\chapter{Materiales y datos}

Este capítulo describe las fuentes de imagen empleadas, la definición de la
tarea de clasificación, la composición del conjunto resultante y el
preprocesado común aplicado antes del modelado. La caracterización de los
datos es central en este trabajo: como se mostrará, una propiedad de la propia
composición del conjunto ---la coincidencia entre la clase clínica y el dataset
de procedencia--- condiciona por completo la interpretación de los resultados
posteriores.

\section{Datasets: BraTS, UPENN-GBM, IXI, NKI Rockland y BTC\_preop}

El trabajo se apoya exclusivamente en datos públicos de resonancia magnética
(RM) cerebral. Para construir el conjunto de entrenamiento y evaluación
principal ---en adelante, el \emph{pool multi-fuente}--- se combinaron cuatro
cohortes: dos aportan estudios con tumor (clase positiva) y dos aportan
estudios de sujetos sin lesión tumoral (clase negativa). De forma separada se
incorporó una quinta cohorte, BTC\_preop, reservada para una evaluación
intra-dominio que se detalla en capítulos posteriores.

\begin{itemize}
  \item \textbf{BraTS 2021} (\emph{Brain Tumor Segmentation Challenge}). Cohorte
  multicéntrica de gliomas con segmentación experta, distribuida a través de
  The Cancer Imaging Archive (TCIA) en el marco del reto
  RSNA--ASNR--MICCAI~\cite{baid2021brats,menze2015brats,bakas2017advancing,clark2013tcia}.
  Aporta estudios de la clase tumor.

  \item \textbf{UPENN-GBM}. Cohorte de glioblastoma de la Universidad de
  Pensilvania, con RM avanzada y datos clínicos, genómicos y
  radiómicos~\cite{bakas2022upenngbm}. Aporta estudios de la clase tumor.

  \item \textbf{IXI} (\emph{Information eXtraction from Images}). Conjunto de RM
  cerebral de sujetos sanos recopilado por el Biomedical Image Analysis Group
  del Imperial College London~\cite{ixi_dataset}. Aporta estudios de la clase
  no-tumor.

  \item \textbf{NKI Rockland}. Muestra de neuroimagen de la iniciativa
  Nathan Kline Institute--Rockland, orientada a acelerar la ciencia del
  descubrimiento en psiquiatría~\cite{nooner2012nki}. Aporta estudios de la
  clase no-tumor.

  \item \textbf{BTC\_preop} (OpenNeuro \texttt{ds001226}). Cohorte preoperatoria
  de tumor cerebral adquirida en un \emph{único} centro (Ghent University
  Hospital), distribuida bajo licencia CC0~\cite{aerts2022btcpreop}. A
  diferencia del pool, contiene ambas clases ---pacientes con tumor y
  controles--- adquiridas en el mismo escáner, por lo que se emplea como cohorte
  intra-dominio. Solo dispone de la modalidad T1 estructural, lo que condiciona
  el diseño del experimento correspondiente (véase la Sección~\ref{sec:modalidades}).
\end{itemize}

La Tabla~\ref{tab:composicion} resume la composición numérica de cada fuente
tras el preprocesado. El pool multi-fuente reúne \textbf{2267} estudios; la
cohorte BTC\_preop se mantiene aparte con \textbf{36} sujetos (\textbf{25} con
tumor y \textbf{11} controles), composición verificada en los resultados de la
auditoría intra-dominio. De los 25 estudios con tumor, según los metadatos del
depósito \texttt{ds001226}, 11 corresponden a gliomas y 14 a
meningiomas~\cite{aerts2022btcpreop}.

\section{Definición de clases y etiquetado}

La tarea se formula como una \textbf{clasificación binaria a nivel de estudio}:
a cada volumen de RM se le asigna una etiqueta $y \in \{0, 1\}$, donde $y=1$
indica presencia de masa tumoral (clase \emph{tumor}) e $y=0$ indica ausencia
de tumor (clase \emph{no-tumor} o sano). No se aborda segmentación ni
localización de la lesión, sino la priorización de estudios con sospecha de
masa, coherente con un escenario de triaje.

El etiquetado se deriva directamente de la procedencia de cada cohorte: BraTS
2021 y UPENN-GBM son colecciones de pacientes oncológicos y se etiquetan en su
totalidad como clase tumor ($y=1$); IXI y NKI Rockland son cohortes de sujetos
sin tumor y se etiquetan como clase no-tumor ($y=0$). En consecuencia,
\emph{cada dataset aporta una única clase}: ninguno contiene simultáneamente
estudios positivos y negativos. El pool multi-fuente resultante contiene
\textbf{1167} estudios positivos y \textbf{1100} negativos (Tabla~\ref{tab:composicion}).

\begin{table}[ht]
  \centering
  \caption[Composición del conjunto de datos]{Composición del conjunto de datos tras el preprocesado. El pool
  multi-fuente agrupa las cuatro primeras cohortes; BTC\_preop se reserva como
  cohorte intra-dominio. Las cifras del pool proceden del resumen de
  preprocesado y la composición de BTC\_preop (36 = 25 tumor / 11 control) de
  los resultados de la auditoría intra-dominio; ambas fuentes son JSON
  versionados.}
  \label{tab:composicion}
  \begin{tabular}{l l c}
    \hline
    \textbf{Dataset} & \textbf{Clase} & \textbf{Estudios ($n$)} \\
    \hline
    BraTS 2021    & Tumor ($y=1$)    & 580 \\
    UPENN-GBM     & Tumor ($y=1$)    & 587 \\
    IXI           & No-tumor ($y=0$) & 577 \\
    NKI Rockland  & No-tumor ($y=0$) & 523 \\
    \hline
    \textbf{Pool multi-fuente} & Tumor + No-tumor & \textbf{2267} (1167 / 1100) \\
    \hline
    BTC\_preop (\texttt{ds001226}) & Tumor + Control & 36 (25 / 11) \\
    \hline
  \end{tabular}
\end{table}

\section{Composición del problema: relación entre clase y dataset de origen}
\label{sec:confound}

La consecuencia directa del esquema de etiquetado anterior es una propiedad
estructural del conjunto que resulta determinante para todo el análisis: la
clase clínica coincide de forma perfecta con el dataset de procedencia. En
términos de conjuntos,
\[
  y = 1 \iff \text{dataset} \in \{\text{BraTS},\ \text{UPENN}\},
  \qquad
  y = 0 \iff \text{dataset} \in \{\text{IXI},\ \text{NKI}\}.
\]
La auditoría de las etiquetas confirma esta correspondencia: cada cohorte
presenta un único valor de etiqueta (BraTS y UPENN únicamente $y=1$; IXI y NKI
únicamente $y=0$). Dicho de otro modo, conocer la fuente de un estudio
determina su clase con certeza, y viceversa.

Esta coincidencia constituye un \emph{confound} dominio$\leftrightarrow$clase: la variable de
interés (presencia de tumor) está completamente confundida con una variable de
molestia (la procedencia del estudio, ligada al equipo de adquisición y al
protocolo de preprocesado de cada cohorte). En presencia de un confound de esta
magnitud, cualquier característica que distinga las cohortes de origen ---por
ejemplo, diferencias sistemáticas de intensidad introducidas por el escáner o
el pipeline--- es suficiente para predecir la etiqueta sin necesidad de
información clínica. Este es precisamente el mecanismo del \emph{shortcut
learning}, por el cual un modelo aprende un atajo espurio en lugar de la señal
relevante~\cite{geirhos2020shortcut}, y es bien conocido que los efectos de
escáner en datos multicéntricos generan este tipo de
atajos~\cite{glocker2019scanner,castro2020causality}.

Es importante subrayar que el confound es una propiedad de la \emph{composición
de los datos}, no un resultado del modelo: existe con independencia de qué
clasificador se entrene. Por esta misma razón, ningún esquema de partición de
los datos del pool puede eliminarlo, y la única vía para evaluar la detección de
tumor sin esta contaminación es disponer de ambas clases adquiridas en el mismo
dominio. Esa es la motivación de incorporar la cohorte intra-dominio BTC\_preop,
cuyo diseño y resultados se desarrollan en los capítulos de diseño experimental
y resultados.

\section{Selección de modalidades T1/T2}
\label{sec:modalidades}

De entre las modalidades estructurales disponibles se seleccionaron las
ponderaciones \textbf{T1} y \textbf{T2}, presentes en las cuatro cohortes del
pool y habituales en la valoración por imagen de tumores cerebrales, donde
aportan información complementaria sobre la masa y el edema
peritumoral~\cite{villanuevameyer2017imaging,martucci2023mri}. El pool
multi-fuente se procesa, por tanto, como entrada de \textbf{dos canales}
(T1 + T2).

La cohorte intra-dominio BTC\_preop constituye la excepción: el depósito
\texttt{ds001226} no incluye T2 estructural, por lo que el experimento
intra-dominio se realiza necesariamente en modalidad \textbf{T1 únicamente}
(un solo canal)~\cite{aerts2022btcpreop}. Salvo por el número de canales de
entrada, el resto del procesamiento es idéntico al del pool, de manera que la
comparación entre regímenes se mantiene lo más homogénea posible. La pérdida de
información que supone prescindir de T2 en ese experimento se discute como
limitación en el capítulo correspondiente.

\section{Análisis exploratorio de los datos}

\subsection{Preprocesado común}

Todas las cohortes se sometieron a un mismo pipeline de preprocesado, de modo
que las diferencias residuales entre datasets no se debieran a tratamientos
distintos. La secuencia, idéntica para el pool y para BTC\_preop, es:

\begin{enumerate}
  \item \textbf{Extracción cerebral} (\emph{skull-stripping}) con HD-BET, una
  herramienta de eliminación automática de cráneo basada en redes
  neuronales~\cite{isensee2019hdbet}.
  \item \textbf{Reorientación} a la convención de ejes RAS.
  \item \textbf{Remuestreo} a una resolución isótropa de \textbf{1\,mm\textsuperscript{3}}.
  \item \textbf{Recorte/relleno} (\emph{crop/pad}) a un volumen común de
  \textbf{192\,$\times$\,224\,$\times$\,192} vóxeles.
  \item \textbf{Normalización} de intensidad mediante \emph{z-score} (media 0,
  desviación típica 1) calculado sobre el tejido cerebral, segmentado con un
  umbral de Otsu.
\end{enumerate}

La resolución y la geometría de salida (1\,mm\textsuperscript{3} e
192\,$\times$\,224\,$\times$\,192 vóxeles) se fijan en el resumen de
preprocesado y son comunes a todas las cohortes, lo que garantiza que los
volúmenes son directamente comparables y apilables como tensores de entrada.

\subsection{Distribución de intensidades por dataset}

El análisis exploratorio se centró en comprobar si las cohortes son distinguibles
a partir de descriptores triviales de intensidad, anteriores a cualquier
modelado. La Figura~\ref{fig:intensity} muestra la distribución de varios
estadísticos de intensidad (fracción de vóxeles no nulos tras el
\emph{skull-stripping}, media y percentiles) agrupados por dataset de origen.

\begin{figure}[ht]
  \centering
  \includegraphics[width=0.95\textwidth]{figuras/intensity_by_dataset.png}
  \caption[Distribuciones de intensidad por dataset]{Distribuciones de estadísticos de intensidad por dataset de origen.
  Las cuatro cohortes del pool son separables ya a partir de descriptores
  triviales (fracción de vóxeles no nulos, media y percentiles de intensidad),
  lo que indica que la firma de procedencia está presente en los píxeles crudos,
  antes de entrar al modelo.}
  \label{fig:intensity}
\end{figure}

Las distribuciones son visiblemente separables por cohorte: incluso descriptores
tan simples como la fracción de vóxeles no nulos ---una huella directa del
\emph{skull-stripping}--- o los percentiles de intensidad difieren de forma
sistemática entre datasets. Dado que, por construcción, la procedencia coincide
con la clase (Sección~\ref{sec:confound}), esta separabilidad en los datos
crudos implica que la etiqueta es predecible a partir de la intensidad, sin
recurrir a información anatómica ni a modelo alguno. Esta observación
exploratoria anticipa y explica los resultados del capítulo de resultados, donde
se cuantifica formalmente la magnitud del confound.
